\DocumentMetadata{
  pdfversion=2.0,
  pdfstandard=ua-2,
  lang=en-US,
  tagging=on,
  tagging-setup={math/setup=mathml-SE,tabsorder=structure}
}
\documentclass[11pt,letterpaper]{article}
\usepackage[margin=0.68in,includefoot]{geometry}
\usepackage{newcomputermodern}
\usepackage{amsmath,amscd,mathtools}
\usepackage{tikz,adjustbox}
\providecommand{\square}{\mdlgwhtsquare}
\usepackage[hidelinks]{hyperref}
\hypersetup{
  pdftitle={Analysis PhD Exam, September 1991},
  pdfauthor={University of Florida Department of Mathematics},
  pdfsubject={Graduate mathematics examination},
  pdfdisplaydoctitle=true,
  bookmarksopen=true
}
\setcounter{secnumdepth}{0}
\setlength{\parindent}{0pt}
\setlength{\parskip}{0.28em}
\setlength{\emergencystretch}{3em}
\setlength{\leftmargini}{2.15em}
\setlength{\leftmarginii}{2.65em}
\setlength{\leftmarginiii}{3em}
\pagestyle{empty}
\raggedbottom
\allowdisplaybreaks
\NewTaggingSocketPlug{block/list/label}{examproblem}{%
  \tagstructbegin{tag=Lbl}\tagstructend%
  \tagstructbegin{tag=\LBody}%
  \tagstructbegin{tag=\ProblemHeadingTag,title-o={\ProblemHeadingText}}%
  \tagmcbegin{tag=\ProblemHeadingTag,actualtext={\ProblemHeadingText}}%
  #2%
  \tagmcend%
  \tagstructend%
}
\newenvironment{examproblems}[2]{%
  \def\ProblemHeadingTag{#2}%
  \AssignTaggingSocketPlug{block/list/label}{examproblem}%
  \begin{enumerate}%
  \setcounter{enumi}{#1}%
  \renewcommand{\labelenumi}{\textbf{\theenumi.}}%
  \setlength{\topsep}{0.35em}%
  \setlength{\partopsep}{0pt}%
  \setlength{\itemsep}{0.48em}%
  \setlength{\parsep}{0.18em}%
}{\end{enumerate}}
\newenvironment{examsubparts}{%
  \AssignTaggingSocketPlug{block/list/label}{default}%
  \begin{enumerate}%
  \setlength{\topsep}{0.18em}%
  \setlength{\partopsep}{0pt}%
  \setlength{\itemsep}{0.16em}%
  \setlength{\parsep}{0.1em}%
}{\end{enumerate}}
\newenvironment{examinstructions}{%
  \par\begingroup\itshape
}{%
  \par\endgroup\vspace{0.18em}
}
\makeatletter
\renewcommand\subsection{\@startsection{subsection}{2}{0pt}%
  {-1.25ex plus -.25ex minus -.1ex}%
  {0.55ex plus .1ex}%
  {\normalfont\large\bfseries}}
\renewcommand{\@maketitle}{%
  \newpage\null\vskip -1em%
  {\footnotesize\bfseries
    \href{https://www.ufl.edu/}{University of Florida}, \href{https://math.ufl.edu/}{Department of Mathematics}\par}%
  \vskip -0.5em%
  \tagstructbegin{tag=H1,title={Analysis PhD Exam, September 1991}}%
  \tagpdfparaOff%
  \tagmcbegin{tag=H1}%
    {\LARGE\bfseries\href{https://gma.math.ufl.edu/exams/analysis-phd/}{Analysis PhD Exam}, September 1991\par}%
  \tagmcend%
  \tagpdfparaOn%
  \tagstructend%
  \vskip 0.25em%
  \tagmcbegin{artifact}%
    \hrule height 1pt%
  \tagmcend%
  \vskip 1em%
}
\makeatother
\title{Analysis PhD Exam, September 1991}
\author{}
\date{}
\begin{document}
\maketitle
\thispagestyle{empty}
\begin{examinstructions}
Present all work in a neat and logical fashion.
\end{examinstructions}
\begin{examproblems}{0}{H2}
\def\ProblemHeadingText{Problem 1.}
\item
State and prove the Lebesgue Monotone Convergence Theorem.
\def\ProblemHeadingText{Problem 2.}
\item
Prove that \(L^1(S,\Sigma,\mu)\) is complete, where~\(\mu\) is a nonnegative measure.
\def\ProblemHeadingText{Problem 3.}
\item
Give the construction of the product measure \(\mu\times\nu\) on \(\mathcal S\times\mathcal W\), where \((X,\mathcal S,\mu)\) and \((Y,\mathcal W,\nu)\) are~\(\sigma\)-finite nonnegative measure spaces.
\def\ProblemHeadingText{Problem 4.}
\item
Let~\(f\) be Lebesgue integrable on \(\mathbb R\). Suppose \(\int_0^x f\,dm=0\) for all real~\(x\). What can you say about~\(f\)? Prove your answer.
\def\ProblemHeadingText{Problem 5.}
\item
Let~\(f\) be a Lebesgue integrable function on \(\mathbb R\). Let \(\epsilon>0\). Is it possible to find an interval~\(I\) such that whenever~\(E\) is a measurable set such that \(E\cap I=\phi\), then \(\left|\int_E f\,dm\right|<\epsilon\)?
\def\ProblemHeadingText{Problem 6.}
\item
Let~\(X\) be a metric space. Let \(\{P_n\}\) be a sequence of regular probability measures on \(\mathcal B(X)\), the Borel subsets of~\(X\). Let \(P=\sum_{n=1}^{\infty}\left(\frac{1}{2^n}\right)P_n\). Show that~\(P\) is regular on \(\mathcal B(X)\).
\def\ProblemHeadingText{Problem 7.}
\item
Let \(\mu(dx)=x^2\,dx\) on \([0,1]\). Let \(T:[0,1]\to[0,1]\) be defined by \(T(x)=x^4\). Compute the Radon–Nikodym derivative of the image measure \(T\mu\).
\def\ProblemHeadingText{Problem 8.}
\item
Let \((\Omega,\mathcal F,\mu)\) be a measure space with \(\mu(\Omega)=1\). Let \(\mathcal A\) and \(\mathcal B\) be sub~\(\sigma\)-algebras of \(\mathcal F\). Suppose \(\mu(A\cap B)=\mu(A)\mu(B)\) for all \(A\in\mathcal A\) and \(B\in\mathcal B\). Prove that \(\int fg\,d\mu=\left(\int f\,d\mu\right)\left(\int g\,d\mu\right)\) for all \(\mathcal A\)-measurable positive functions~\(f\) and all \(\mathcal B\)-measurable positive functions~\(g\).
\end{examproblems}
\end{document}
